\chapter{Crisis Doxonomics}
\label{ch:crisis-doxonomics}

\epigraph{In a crisis, the old world is dying and the new world struggles to be born. Now is the time of monsters.}{Antonio Gramsci (paraphrased)}

\noindent
Crises---wars, pandemics, economic collapses, ecological disasters---are periods of accelerated belief change.
The normal inertia of common sense is temporarily suspended.
This chapter analyzes the doxonomics of crisis: how fear amplifies narrative capture, how emergencies expand the Overton window, and how the aftermath determines which new common sense is installed.

\section{Crisis as Phase Transition}
\label{sec:ch29-phase}

\begin{model}[Crisis Belief Dynamics]
\label{mod:crisis}
\index{crisis doxonomics}
In normal times, belief changes slowly:
\[
  \frac{d\beliefvec}{dt} = \varepsilon \cdot f(\fundingvec, \beliefvec) \qquad \varepsilon \ll 1
\]
During a crisis, the time constant accelerates:
\[
  \frac{d\beliefvec}{dt} = f(\fundingvec, \beliefvec) \qquad (\varepsilon \to 1)
\]
The crisis does not change the structure of doxonomic dynamics; it changes their \emph{speed}.
Beliefs that would take decades to shift can shift in months.
\end{model}

\section{Fear as Amplifier}
\label{sec:ch29-fear}

\begin{proposition}[Fear and Credibility Weighting]
\label{prop:fear}
\index{fear}
Under conditions of threat, agents increase the credibility weight assigned to authority sources:
\[
  \credibility{k}^{\text{crisis}} = \credibility{k} + \eta \cdot \text{threat}(t)
\]
for institutional channels (government, experts, legacy media), while decreasing weights for unfamiliar or non-authoritative sources.
This concentrates doxonomic power in established institutions during crises.
\end{proposition}

\section{Overton Window Shifts}
\label{sec:ch29-overton}

\begin{definition}[Overton Window]
\label{def:overton}
\index{Overton window}
The \emph{Overton window} is the set of propositions $W \subset \narspace$ that are perceived as politically viable at time $t$.
Propositions outside $W$ are ``unthinkable'' regardless of their merit.
A crisis can shift $W$ dramatically, making previously unthinkable policies (surveillance, nationalization, border closure, universal basic income) temporarily thinkable.
\end{definition}

The doxonomic contest during a crisis is over which direction $W$ shifts.
The group that captures the post-crisis narrative determines the new Overton window.

\section{Notes and References}
\label{sec:ch29-notes}

On crisis and belief change, see \textcite{klein2007}.
Gramsci's concept of interregnum connects to \textcite{gramsci1971}.
The Overton window concept is from Joseph Overton (Mackinac Center, 2003).
On emergency governance and narrative, see \textcite{lippmann1922} and \textcite{herman1988}.

\section{Exercises}

\begin{exercise}
Using Model~\ref{mod:crisis}, analyze the COVID-19 pandemic as a doxonomic event.
Which beliefs shifted rapidly?
Which institutions captured the narrative?
What new common sense was installed?
\end{exercise}

\begin{exercise}
Estimate the Overton window before and after the 2008 financial crisis.
Which propositions entered or exited the window?
\end{exercise}

\begin{exercise}
Design a ``doxonomic preparedness'' framework for progressive movements.
What narratives should be ready for deployment when the next crisis opens the Overton window?
\end{exercise}

\begin{exercise}
Model two competing groups racing to capture the post-crisis narrative.
Formalize as a game and identify the Nash equilibrium investment levels.
\end{exercise}
